\chapter{Communication Complexity — Public Coin Approximation}
%%% generated by the blueprint pipeline from TCSlib.CommunicationComplexity.NewmanTheorem.PublicCoinApproximation
\section{Overview}

This module provides the bridge between public-coin protocols over arbitrary finite
probability spaces and those using the canonical \texttt{CoinTape} randomness source.
The key result is that any finite-message public-coin protocol can be approximated by a
\texttt{CommunicationComplexity.CoinTape}-based protocol of equal complexity, with the approximation error
degrading by at most the chosen slack $\delta > 0$.

%%%%%%%%%%%%%%%%%%%%%%%%%%%%%%%%%%%%%%%%%%%%%%%%%%%%%%%%%%%%%%%%%%%%%%%%%%%%%%%
\section{Declarations}

\begin{definition}[Conversion of a public-coin protocol to CoinTape]
\lean{CommunicationComplexity.PublicCoin.FiniteMessage.Protocol.toCoinTape}
\label{CommunicationComplexity.PublicCoin.FiniteMessage.Protocol.toCoinTape}
\leanok
\uses{CommunicationComplexity.CoinTape, CommunicationComplexity.Deterministic.FiniteMessage.Protocol.comap, CommunicationComplexity.FiniteProbabilitySpace, CommunicationComplexity.Internal.single_coin_approx, CommunicationComplexity.PublicCoin.FiniteMessage.Protocol}
Given a public-coin finite-message protocol $p$ over an arbitrary finite probability
space $\Omega$ and a slack $\delta > 0$, this definition produces a natural number $n$
together with a finite-message protocol over $\mathrm{CoinTape}(n)$ that approximates
$p$. The construction works by pulling back $p$ along a measure-approximating map
$\varphi : \mathrm{CoinTape}(n) \to \Omega$ provided by
\texttt{CommunicationComplexity.Internal.single\_coin\_approx}.
\end{definition}

\begin{theorem}[Coin-tape conversion preserves complexity]
\lean{CommunicationComplexity.PublicCoin.FiniteMessage.Protocol.toCoinTape_complexity}
\label{CommunicationComplexity.PublicCoin.FiniteMessage.Protocol.toCoinTape_complexity}
\leanok
\uses{CommunicationComplexity.Deterministic.FiniteMessage.Protocol.complexity, CommunicationComplexity.FiniteProbabilitySpace, CommunicationComplexity.PublicCoin.FiniteMessage.Protocol, CommunicationComplexity.PublicCoin.FiniteMessage.Protocol.toCoinTape}
\difficulty{1}
Let $\Omega$ be a finite probability space, and let $p$ be a public-coin finite-message
protocol over shared randomness $\Omega$ with input sets $X$, $Y$ and output type
$\alpha$. For every real number $\delta > 0$, the coin-tape conversion of $p$ at slack
$\delta$ yields a length $n$ and a finite-message protocol over $\mathrm{CoinTape}(n)$
whose communication complexity equals the communication complexity of $p$.
\end{theorem}

\begin{theorem}[Coin-tape conversion preserves approximate satisfaction]
\lean{CommunicationComplexity.PublicCoin.FiniteMessage.Protocol.toCoinTape_approxSatisfies}
\label{CommunicationComplexity.PublicCoin.FiniteMessage.Protocol.toCoinTape_approxSatisfies}
\leanok
\uses{CommunicationComplexity.CoinTape, CommunicationComplexity.Deterministic.FiniteMessage.Protocol.comap, CommunicationComplexity.Deterministic.FiniteMessage.Protocol.comap_run, CommunicationComplexity.FiniteProbabilitySpace, CommunicationComplexity.Internal.single_coin_approx, CommunicationComplexity.PublicCoin.FiniteMessage.Protocol, CommunicationComplexity.PublicCoin.FiniteMessage.Protocol.ApproxSatisfies, CommunicationComplexity.PublicCoin.FiniteMessage.Protocol.rrun, CommunicationComplexity.PublicCoin.FiniteMessage.Protocol.toCoinTape}
\difficulty{5}
Let $p$ be a public-coin finite-message protocol over a finite probability space
$\Omega$, with private input sets $X$ and $Y$ and outputs in $\alpha$; let $Q : X \to Y
\to \alpha \to \mathrm{Prop}$ be a predicate; and let $\varepsilon$ and $\delta$ be real
numbers with $\delta > 0$. If $p$ approximately satisfies $Q$ with error at most
$\varepsilon$, then the coin-tape protocol obtained by converting $p$ with slack
$\delta$ approximately satisfies $Q$ with error at most $\varepsilon + \delta$.
\end{theorem}
